\documentclass[aspectratio=169]{beamer}
% Shared Beamer configuration for MUS2640 lecture decks (included after \documentclass).
\usepackage[utf8]{inputenc}
\usepackage[T1]{fontenc}
\usepackage{lmodern}
\usepackage{graphicx}
\usepackage{booktabs}

\usetheme{Madrid}
\usecolortheme{dolphin}
\usefonttheme{professionalfonts}

\setbeamercovered{transparent}
\setbeamertemplate{navigation symbols}{}
\setbeamertemplate{footline}[frame number]

\newcommand{\coursename}{MUS2640 · Sensing Sound and Music}
\newcommand{\instituteuo}{University of Oslo · Department of Musicology / RITMO}


\title{Week 9 · Physiology}
\subtitle{\coursename}
\institute{\instituteuo}
\date{}

\begin{document}

\begin{frame}
  \titlepage
\end{frame}

\begin{frame}{Aims}
  \begin{itemize}
    \item Link affect and arousal to \textbf{autonomic} signals: HR, HRV, EDA, respiration.
    \item Compare performer vs listener physiology; interpersonal synchrony (preview).
    \item Measurement basics — what each signal can and cannot imply.
  \end{itemize}
\end{frame}

\begin{frame}{Roadmap}
  \begin{enumerate}
    \item Sympathetic vs parasympathetic balance (intuition).
    \item Signals: ECG-derived HR/HRV; skin conductance; breathing; muscle tension (categories).
    \item Laboratory vs concert / wearable constraints.
    \item Ethics: consent, context, interpretation humility.
  \end{enumerate}
\end{frame}

\begin{frame}{Bridge from embodiment}
  \begin{itemize}
    \item Same musical event can ``move'' people differently — distributions + context matter.
    \item Self-report + physiology + behaviour often triangulate (not replace each other).
  \end{itemize}
\end{frame}

\begin{frame}{In-class}
  \begin{itemize}
    \item Discuss one study design: what would you measure during a live concert?
    \item Relate one physiological idea to a personal listening memory (short round).
  \end{itemize}
\end{frame}

\begin{frame}{Outside class}
  \begin{itemize}
    \item Chapter exercises; connect labels (HRV, tonic/phasic EDA) to plain-language descriptions.
  \end{itemize}
\end{frame}

\begin{frame}{Summary}
  \begin{itemize}
    \item Physiology grounds ``felt experience'' in measurable responses — with caveats.
    \item Next: \textbf{Vision} — AV integration, attention, eye-tracking in performance research.
  \end{itemize}
\end{frame}

\end{document}
